\chapter{Conclusion\label{cha:chapter7}}


\section{Summary\label{sec:summary}}

Previous research has shown that sparse or coarse satellite SIF can be reconstructed or spatially enhanced using optical, radiation, meteorological, land-cover, and spatio-temporal predictors \cite{li2019global,zhang2018global,shen2022spatiotemporal}. Global reconstruction products improve spatial and temporal continuity, but their fine-scale values are inferred mainly from predictor relationships and are not generally required to preserve a contemporaneous coarse SIF observation. TROPOMI-based downscaling methods can retain more of the measured signal through bias correction, redistribution, or physical conservation constraints, although they depend on the availability of a parent observation for the location and date being enhanced \cite{hu2024spatially,chen2025improved,zhang2023generating}. CNN-based studies further demonstrate that spatial predictor grids can represent neighbourhood patterns and heterogeneous land cover that are unavailable to pixel-wise tabular models \cite{gensheimer2022convolutional,du2025cnsif}. However, a finer output grid does not ensure that the estimated subpixel pattern represents true physiological variation, because optical predictors can describe canopy structure more readily than short-term photosynthetic regulation. Most widely evaluated enhanced products also remain at approximately 500~m resolution or coarser, which is insufficient for separating many small and irregularly shaped agricultural fields. This limitation is more difficult to address with OCO-2 than with TROPOMI because the narrow and widely separated OCO-2 tracks leave large areas without a contemporaneous SIF value that could constrain a redistribution method \cite{yu2019high}. Validation is another common limitation, as many products are assessed against held-out supervisory observations or indirect proxies such as GPP and yield rather than independent SIF measured at the stated fine resolution. Consequently, agreement after aggregation to a satellite footprint does not establish the accuracy of every fine-grid pixel. Crop-yield studies nevertheless indicate that SIF contains useful information about crop productivity and stress, although its advantage depends on crop type, observation timing, geographic scale, and the inclusion of weather or other environmental predictors \cite{peng2020assessing,song2018satellite,qiu2022monitoring}. These gaps motivated the present study to test OCO-2 enhancement with Sentinel-2 predictors at 20~m while explicitly separating output resolution from validation support and examining whether the resulting crop-pure signal benefits winter-wheat yield prediction.

The experimental workflow consisted of two connected stages: SIF enhancement followed by regional winter-wheat yield prediction. The enhancement stage used quality-controlled OCO-2 observations from February to July during 2019--2024 across nine German Sentinel-2 MGRS tiles, with two further tiles reserved for external evaluation. SIF retrieved near 757 and 771~nm was combined into a noise-reduced target and matched to Sentinel-2 reflectance and spectral indices, GLASS vegetation variables, VIIRS radiation data, seasonal encodings, and annual crop-composition information. All raster predictors were aligned to a common 20~m grid, while the irregular OCO-2 polygons were retained as fractional masks so that predictions could be averaged over their actual spatial support. Several model and data representations were evaluated, including polygon-mean Random Forest and XGBoost models, an area-to-point multilayer perceptron, single- and multi-footprint convolutional models, and a U-Net trained on fixed 4~km spatial aggregates. These experiments tested both whether a spatial model could produce context-aware fine-grid maps and whether OCO-2 aggregation could provide a more stable training target. The resulting 20~m surfaces were treated as model-derived downscaled estimates rather than direct SIF observations. This design allowed the effects of target construction, predictor representation, spatial support, and model family to be examined within one workflow, although the experiments were not a controlled comparison of architecture alone.

Training, validation, and testing partitions were constructed so that complete dates, footprints, or connected date-product components remained together according to the requirements of each experiment. The principal model averaged at least five same-date and same-mode OCO-2 observations within each fixed 4~km cell and trained a U-Net whose output was aggregated through the corresponding footprint-weight map. Internal evaluation used held-out aggregate cells, whereas transfer was assessed by applying the saved model and validation-derived calibration to native OCO-2 footprints in two previously unseen MGRS tiles. Performance was additionally examined by month, year, measurement mode, plant-hardiness zone, SIF interval, tile, and footprint count to identify conditions associated with unstable predictions. For the agricultural stage, the selected model was applied to winter-wheat pixels in five Bavarian MGRS tiles and summarized from March to July for each NUTS-3 region and year. Random Forest yield models were trained on 2019--2022 and tested on 2024 using either crop-pure enhanced SIF, raw footprints containing at least 5\% winter wheat, or all accepted raw footprints. The response data and model configuration were held constant so that the comparison primarily reflected the alternative SIF representations. Paired bootstrap resampling of the 48 test regions was then used to quantify uncertainty in the differences between the crop-pure and raw-SIF models.

The principal 4~km spatial-aggregate U-Net achieved the strongest internal result, with an RMSE of 0.0701 and an $R^2$ of 0.761 on 931 held-out grid cells. Aggregation reduced both target and residual variability, confirming that averaging several OCO-2 observations can provide more stable supervision, but it also narrowed the SIF range and removed most extreme values. When transferred to individual footprints in two unseen tiles, performance decreased to an RMSE of 0.1701 and an $R^2$ of 0.373, showing that internal aggregate accuracy overstates performance at noisier native support. Native-footprint U-Net and tree-based models produced broadly similar aggregate metrics, indicating that convolutional architecture alone did not guarantee greater predictive accuracy than carefully constructed tabular predictors. The convolutional models nevertheless retained an important representational advantage because their 20~m maps used neighbouring pixels and preserved field boundaries and landscape context. Errors generally increased later in the growing season, when observed SIF values and their variability were also greater, while sparsely sampled zone-month combinations remained too uncertain for strong agro-climatic conclusions. In the yield experiment, crop-pure enhanced SIF produced the lowest 2024 RMSE at 0.9406~t/ha, compared with 1.0285~t/ha for the 5\%-wheat raw-footprint model and 1.0874~t/ha when all raw footprints were used. This corresponds to RMSE reductions of approximately 8.5\% and 13.5\%, supporting the general literature finding that spatially targeted SIF can provide useful crop information \cite{peng2020assessing,song2018satellite,qiu2022monitoring}. However, the crop-pure model still had a negative $R^2$ of $-0.1456$, and both 95\% bootstrap intervals for its RMSE improvement crossed zero, so the apparent advantage was not statistically conclusive. The results therefore agree with SIFnet and CNSIF in showing the value of spatial predictor structure for heterogeneous landscapes, but they also demonstrate that fine output resolution must not be confused with fine-scale validation \cite{gensheimer2022convolutional,du2025cnsif}. Unlike observation-preserving TROPOMI methods, the present OCO-2 framework can predict between satellite tracks without a contemporaneous parent measurement, but this flexibility increases reliance on the auxiliary predictors and makes independent fine-resolution validation a central requirement for future work \cite{chen2025improved,zhang2023generating}.

\section{Dissemination\label{sec:dissemination}}

The principal outputs of this thesis are the evaluated SIF-enhancement approach, the associated data-preparation and modelling workflow, and evidence concerning the usefulness and limitations of crop-pure SIF for regional yield prediction. The most immediate potential users are remote-sensing, SIF, and machine-learning researchers who require methods for combining sparse fluorescence observations with spatially complete high-resolution predictors \cite{li2019global,gensheimer2022convolutional,du2025cnsif}. Agricultural researchers and crop-monitoring practitioners could use the resulting crop-specific SIF summaries together with weather, soil, and management information to investigate crop productivity, stress, and yield. More broadly, spatially detailed SIF estimates could support research on GPP, vegetation phenology, and plant responses to drought or heat, applications that have also motivated previous SIF products and agricultural studies \cite{qiu2022monitoring,song2018satellite,peng2020assessing}. The thesis itself is the main formal dissemination output and documents the datasets, processing decisions, model configurations, validation design, and interpretation limits needed to understand and reproduce the work. The workflow should currently be regarded as a research prototype, and no separate peer-reviewed publication or integration into an operational monitoring service is claimed. Integration into a larger environment would mean connecting its data acquisition, preprocessing, model inference, uncertainty reporting, and map outputs to an established Earth-observation, GIS, agricultural-monitoring, or decision-support platform. Such adoption would require more automated data processing, sufficient computing infrastructure, stable data interfaces, independent fine-resolution validation, and continued performance monitoring before the enhanced SIF maps could be used operationally.

\section{Problems Encountered\label{sec:problems}}

Summarize the main problems. How did you solve them? Why didn't you solve them?

\section{Outlook\label{sec:outlook}}

Future work will enhance Component X with new services and features that can be used ...
